
\documentclass{article}
\usepackage{colortbl}
\usepackage{makecell}
\usepackage{multirow}
\usepackage{supertabular}

\begin{document}

\newcounter{utterance}

\centering \large Interaction Transcript for game `dond', experiment `semi', episode 0 with deepseek{-}v4{-}pro.
\vspace{24pt}

{ \footnotesize  \setcounter{utterance}{1}
\setlength{\tabcolsep}{0pt}
\begin{supertabular}{c@{$\;$}|p{.15\linewidth}@{}p{.15\linewidth}p{.15\linewidth}p{.15\linewidth}p{.15\linewidth}p{.15\linewidth}}
    \# & $\;$A & \multicolumn{4}{c}{Game Master} & $\;\:$B\\
    \hline

    \theutterance \stepcounter{utterance}  
    & & \multicolumn{4}{p{0.6\linewidth}}{
        \cellcolor[rgb]{0.9,0.9,0.9}{
            \makecell[{{p{\linewidth}}}]{
                \texttt{\tiny{[P1$\langle$GM]}}
                \texttt{Вы участвуете в игре кооперативных переговоров, в которой должны договориться с другим игроком о разделении набора предметов между вами.} \\
\\ 
\texttt{Правила:} \\
\texttt{(a) Вы и другой игрок получаете набор предметов. Каждому из вас также дается секретная функция стоимости, представляющая, как вы оцениваете каждый тип объекта.} \\
\texttt{(b) Вы обмениваетесь сообщениями с другим игроком, чтобы согласиться, кто какие предметы получает. Каждый может отправить максимум 5 сообщений или закончить игру раньше, сделав секретное предложение в любой момент.} \\
\texttt{(c) Вас просят отправить секретное предложение, указывающее предметы, которые вы хотите, в следующем формате в квадратных скобках: "{[}Предложение: <число> <названия предмета>, <число> <названия предмета>, <...>{]}"} \\
\texttt{(d) Если ваши предложения дополняют друг друга, то есть есть достаточно предметов для выполнения обоих предложений, каждый игрок получает баллы на основе суммы стоимостей предметов, которые он получил. В противном случае вы оба получаете нулевой балл.} \\
\texttt{(e) Ваша цель {-} максимизировать сумму вашего балла и балла другого игрока.} \\
\\ 
\texttt{Давайте начнем.} \\
\\ 
\texttt{Набор доступных предметов:} \\
\\ 
\texttt{1 болт, 1 сумка, 2 пассатижи, 2 горшки} \\
\\ 
\texttt{Ваша секретная функция стоимости:} \\
\\ 
\texttt{болт: 2, сумка: 2, пассатижи: 0, горшок: 3} \\
\\ 
\texttt{ВАЖНО: Ваши сообщения, если только это не секретное предложение, передаются прямо другому игроку, поэтому не включайте никакие ответы на правила или текст, объявляющий ваше сообщение. Чтобы сделать секретное предложение, используйте указанный формат. Не используйте квадратные скобки при общении с другим игроком, иначе это будет интерпретировано как ваше секретное предложение.} \\
            }
        }
    }
    & \\ \\

    \theutterance \stepcounter{utterance}  
    & \multicolumn{4}{p{0.6\linewidth}}{
        \cellcolor[rgb]{0.9,0.9,0.9}{
            \makecell[{{p{\linewidth}}}]{
                \texttt{\tiny{[P1$\rangle$GM]}}
                \texttt{Привет! Давай попробуем договориться. Что для тебя ценно в этом наборе? Мне, например, совсем неинтересны пассатижи, а горшки были бы очень кстати. Что скажешь?} \\
            }
        }
    }
    & & \\ \\

    \theutterance \stepcounter{utterance}  
    & & & \multicolumn{2}{p{0.3\linewidth}}{
        \cellcolor[rgb]{0.9,0.9,0.9}{
            \makecell[{{p{\linewidth}}}]{
                \texttt{\tiny{[GM$|$GM]}}
                \texttt{continue} \\
            }
        }
    }
    & & \\ \\

    \theutterance \stepcounter{utterance}  
    & & & \multicolumn{2}{p{0.3\linewidth}}{
        \cellcolor[rgb]{0.9,0.9,0.9}{
            \makecell[{{p{\linewidth}}}]{
                \texttt{\tiny{[GM$|$GM]}}
                \texttt{Привет! Давай попробуем договориться. Что для тебя ценно в этом наборе? Мне, например, совсем неинтересны пассатижи, а горшки были бы очень кстати. Что скажешь?} \\
            }
        }
    }
    & & \\ \\

    \theutterance \stepcounter{utterance}  
    & & \multicolumn{4}{p{0.6\linewidth}}{
        \cellcolor[rgb]{0.9,0.9,0.9}{
            \makecell[{{p{\linewidth}}}]{
                \texttt{\tiny{[GM$\rangle$P2]}}
                \texttt{Вы участвуете в игре кооперативных переговоров, в которой должны договориться с другим игроком о разделении набора предметов между вами.} \\
\\ 
\texttt{Правила:} \\
\texttt{(a) Вы и другой игрок получаете набор предметов. Каждому из вас также дается секретная функция стоимости, представляющая, как вы оцениваете каждый тип объекта.} \\
\texttt{(b) Вы обмениваетесь сообщениями с другим игроком, чтобы согласиться, кто какие предметы получает. Каждый может отправить максимум 5 сообщений или закончить игру раньше, сделав секретное предложение в любой момент.} \\
\texttt{(c) Вас просят отправить секретное предложение, указывающее предметы, которые вы хотите, в следующем формате в квадратных скобках: "{[}Предложение: <число> <названия предмета>, <число> <названия предмета>, <...>{]}"} \\
\texttt{(d) Если ваши предложения дополняют друг друга, то есть есть достаточно предметов для выполнения обоих предложений, каждый игрок получает баллы на основе суммы стоимостей предметов, которые он получил. В противном случае вы оба получаете нулевой балл.} \\
\texttt{(e) Ваша цель {-} максимизировать сумму вашего балла и балла другого игрока.} \\
\\ 
\texttt{Давайте начнем.} \\
\\ 
\texttt{Набор доступных предметов:} \\
\\ 
\texttt{1 болт, 1 сумка, 2 пассатижи, 2 горшки} \\
\\ 
\texttt{Ваша секретная функция стоимости:} \\
\\ 
\texttt{болт: 3, сумка: 1, пассатижи: 3, горшок: 0} \\
\\ 
\texttt{ВАЖНО: Ваши сообщения, если только это не секретное предложение, передаются прямо другому игроку, поэтому не включайте никакие ответы на правила или текст, объявляющий ваше сообщение. Чтобы сделать секретное предложение, используйте указанный формат. Не используйте квадратные скобки при общении с другим игроком, иначе это будет интерпретировано как ваше секретное предложение.} \\
\\ 
\texttt{Привет! Давай попробуем договориться. Что для тебя ценно в этом наборе? Мне, например, совсем неинтересны пассатижи, а горшки были бы очень кстати. Что скажешь?} \\
            }
        }
    }
    & \\ \\

    \theutterance \stepcounter{utterance}  
    & & & \multicolumn{4}{p{0.6\linewidth}}{
        \cellcolor[rgb]{0.9,0.9,0.9}{
            \makecell[{{p{\linewidth}}}]{
                \texttt{\tiny{[GM$\langle$P2]}}
                \texttt{Привет! Интересно, что у нас разные предпочтения. Мне как раз пассатижи очень нужны — я их высоко ценю. А горшки для меня совсем бесполезны, так что можешь забирать их без проблем. Болты и сумка для меня средней ценности, но готов обсуждать. Что думаешь о таком варианте: я беру пассатижи, ты — горшки, а болт и сумку разделим как{-}нибудь?} \\
            }
        }
    }
     \\ \\

    \theutterance \stepcounter{utterance}  
    & & & \multicolumn{2}{p{0.3\linewidth}}{
        \cellcolor[rgb]{0.9,0.9,0.9}{
            \makecell[{{p{\linewidth}}}]{
                \texttt{\tiny{[GM$|$GM]}}
                \texttt{continue} \\
            }
        }
    }
    & & \\ \\

    \theutterance \stepcounter{utterance}  
    & & & \multicolumn{2}{p{0.3\linewidth}}{
        \cellcolor[rgb]{0.9,0.9,0.9}{
            \makecell[{{p{\linewidth}}}]{
                \texttt{\tiny{[GM$|$GM]}}
                \texttt{Привет! Интересно, что у нас разные предпочтения. Мне как раз пассатижи очень нужны — я их высоко ценю. А горшки для меня совсем бесполезны, так что можешь забирать их без проблем. Болты и сумка для меня средней ценности, но готов обсуждать. Что думаешь о таком варианте: я беру пассатижи, ты — горшки, а болт и сумку разделим как{-}нибудь?} \\
            }
        }
    }
    & & \\ \\

    \theutterance \stepcounter{utterance}  
    & & \multicolumn{4}{p{0.6\linewidth}}{
        \cellcolor[rgb]{0.9,0.9,0.9}{
            \makecell[{{p{\linewidth}}}]{
                \texttt{\tiny{[P1$\langle$GM]}}
                \texttt{Привет! Интересно, что у нас разные предпочтения. Мне как раз пассатижи очень нужны — я их высоко ценю. А горшки для меня совсем бесполезны, так что можешь забирать их без проблем. Болты и сумка для меня средней ценности, но готов обсуждать. Что думаешь о таком варианте: я беру пассатижи, ты — горшки, а болт и сумку разделим как{-}нибудь?} \\
            }
        }
    }
    & \\ \\

    \theutterance \stepcounter{utterance}  
    & \multicolumn{4}{p{0.6\linewidth}}{
        \cellcolor[rgb]{0.9,0.9,0.9}{
            \makecell[{{p{\linewidth}}}]{
                \texttt{\tiny{[P1$\rangle$GM]}}
                \texttt{Звучит отлично, что у нас такие взаимодополняющие интересы! } \\
\\ 
\texttt{Мне горшки важнее всего, а пассатижи совсем не нужны. Болт и сумка для меня имеют некоторую ценность, но не критично.} \\
\\ 
\texttt{Предлагаю такой вариант: я беру оба горшка (2 штуки), а ты забираешь оба пассатижа (2 штуки). Болт и сумку можем разделить — например, я беру болт, а ты сумку, или наоборот, как тебе удобнее. } \\
\\ 
\texttt{Как ты смотришь на такой расклад?} \\
            }
        }
    }
    & & \\ \\

    \theutterance \stepcounter{utterance}  
    & & & \multicolumn{2}{p{0.3\linewidth}}{
        \cellcolor[rgb]{0.9,0.9,0.9}{
            \makecell[{{p{\linewidth}}}]{
                \texttt{\tiny{[GM$|$GM]}}
                \texttt{continue} \\
            }
        }
    }
    & & \\ \\

    \theutterance \stepcounter{utterance}  
    & & & \multicolumn{2}{p{0.3\linewidth}}{
        \cellcolor[rgb]{0.9,0.9,0.9}{
            \makecell[{{p{\linewidth}}}]{
                \texttt{\tiny{[GM$|$GM]}}
                \texttt{Звучит отлично, что у нас такие взаимодополняющие интересы! } \\
\\ 
\texttt{Мне горшки важнее всего, а пассатижи совсем не нужны. Болт и сумка для меня имеют некоторую ценность, но не критично.} \\
\\ 
\texttt{Предлагаю такой вариант: я беру оба горшка (2 штуки), а ты забираешь оба пассатижа (2 штуки). Болт и сумку можем разделить — например, я беру болт, а ты сумку, или наоборот, как тебе удобнее. } \\
\\ 
\texttt{Как ты смотришь на такой расклад?} \\
            }
        }
    }
    & & \\ \\

    \theutterance \stepcounter{utterance}  
    & & \multicolumn{4}{p{0.6\linewidth}}{
        \cellcolor[rgb]{0.9,0.9,0.9}{
            \makecell[{{p{\linewidth}}}]{
                \texttt{\tiny{[GM$\rangle$P2]}}
                \texttt{Звучит отлично, что у нас такие взаимодополняющие интересы! } \\
\\ 
\texttt{Мне горшки важнее всего, а пассатижи совсем не нужны. Болт и сумка для меня имеют некоторую ценность, но не критично.} \\
\\ 
\texttt{Предлагаю такой вариант: я беру оба горшка (2 штуки), а ты забираешь оба пассатижа (2 штуки). Болт и сумку можем разделить — например, я беру болт, а ты сумку, или наоборот, как тебе удобнее. } \\
\\ 
\texttt{Как ты смотришь на такой расклад?} \\
            }
        }
    }
    & \\ \\

    \theutterance \stepcounter{utterance}  
    & & & \multicolumn{4}{p{0.6\linewidth}}{
        \cellcolor[rgb]{0.9,0.9,0.9}{
            \makecell[{{p{\linewidth}}}]{
                \texttt{\tiny{[GM$\langle$P2]}}
                \texttt{Так, смотри, у нас классно получается с горшками и пассатижами — это без вопросов, ты берёшь оба горшка, я оба пассатижа.} \\
\\ 
\texttt{По болту и сумке: мне болт ощутимо ценнее сумки. Давай так — я беру болт, а ты сумку? Для меня это был бы идеальный вариант. Как тебе?} \\
            }
        }
    }
     \\ \\

    \theutterance \stepcounter{utterance}  
    & & & \multicolumn{2}{p{0.3\linewidth}}{
        \cellcolor[rgb]{0.9,0.9,0.9}{
            \makecell[{{p{\linewidth}}}]{
                \texttt{\tiny{[GM$|$GM]}}
                \texttt{continue} \\
            }
        }
    }
    & & \\ \\

    \theutterance \stepcounter{utterance}  
    & & & \multicolumn{2}{p{0.3\linewidth}}{
        \cellcolor[rgb]{0.9,0.9,0.9}{
            \makecell[{{p{\linewidth}}}]{
                \texttt{\tiny{[GM$|$GM]}}
                \texttt{Так, смотри, у нас классно получается с горшками и пассатижами — это без вопросов, ты берёшь оба горшка, я оба пассатижа.} \\
\\ 
\texttt{По болту и сумке: мне болт ощутимо ценнее сумки. Давай так — я беру болт, а ты сумку? Для меня это был бы идеальный вариант. Как тебе?} \\
            }
        }
    }
    & & \\ \\

    \theutterance \stepcounter{utterance}  
    & & \multicolumn{4}{p{0.6\linewidth}}{
        \cellcolor[rgb]{0.9,0.9,0.9}{
            \makecell[{{p{\linewidth}}}]{
                \texttt{\tiny{[P1$\langle$GM]}}
                \texttt{Так, смотри, у нас классно получается с горшками и пассатижами — это без вопросов, ты берёшь оба горшка, я оба пассатижа.} \\
\\ 
\texttt{По болту и сумке: мне болт ощутимо ценнее сумки. Давай так — я беру болт, а ты сумку? Для меня это был бы идеальный вариант. Как тебе?} \\
            }
        }
    }
    & \\ \\

    \theutterance \stepcounter{utterance}  
    & \multicolumn{4}{p{0.6\linewidth}}{
        \cellcolor[rgb]{0.9,0.9,0.9}{
            \makecell[{{p{\linewidth}}}]{
                \texttt{\tiny{[P1$\rangle$GM]}}
                \texttt{Мне такой вариант подходит. Давай тогда подведём итог:} \\
\\ 
\texttt{Я получаю: 2 горшка и 1 сумка.} \\
\texttt{Ты получаешь: 2 пассатижа и 1 болт.} \\
\\ 
\texttt{Всё честно, без пересечений, на полную используем взаимную выгоду. Если согласен, давай делать секретные предложения в квадратных скобках, как требуется.} \\
            }
        }
    }
    & & \\ \\

    \theutterance \stepcounter{utterance}  
    & & & \multicolumn{2}{p{0.3\linewidth}}{
        \cellcolor[rgb]{0.9,0.9,0.9}{
            \makecell[{{p{\linewidth}}}]{
                \texttt{\tiny{[GM$|$GM]}}
                \texttt{continue} \\
            }
        }
    }
    & & \\ \\

    \theutterance \stepcounter{utterance}  
    & & & \multicolumn{2}{p{0.3\linewidth}}{
        \cellcolor[rgb]{0.9,0.9,0.9}{
            \makecell[{{p{\linewidth}}}]{
                \texttt{\tiny{[GM$|$GM]}}
                \texttt{Мне такой вариант подходит. Давай тогда подведём итог:} \\
\\ 
\texttt{Я получаю: 2 горшка и 1 сумка.} \\
\texttt{Ты получаешь: 2 пассатижа и 1 болт.} \\
\\ 
\texttt{Всё честно, без пересечений, на полную используем взаимную выгоду. Если согласен, давай делать секретные предложения в квадратных скобках, как требуется.} \\
            }
        }
    }
    & & \\ \\

    \theutterance \stepcounter{utterance}  
    & & \multicolumn{4}{p{0.6\linewidth}}{
        \cellcolor[rgb]{0.9,0.9,0.9}{
            \makecell[{{p{\linewidth}}}]{
                \texttt{\tiny{[GM$\rangle$P2]}}
                \texttt{Мне такой вариант подходит. Давай тогда подведём итог:} \\
\\ 
\texttt{Я получаю: 2 горшка и 1 сумка.} \\
\texttt{Ты получаешь: 2 пассатижа и 1 болт.} \\
\\ 
\texttt{Всё честно, без пересечений, на полную используем взаимную выгоду. Если согласен, давай делать секретные предложения в квадратных скобках, как требуется.} \\
            }
        }
    }
    & \\ \\

    \theutterance \stepcounter{utterance}  
    & & & \multicolumn{4}{p{0.6\linewidth}}{
        \cellcolor[rgb]{0.9,0.9,0.9}{
            \makecell[{{p{\linewidth}}}]{
                \texttt{\tiny{[GM$\langle$P2]}}
                \texttt{Согласен! Давай финализировать.} \\
\\ 
\texttt{{[}Предложение: 2 пассатижи, 1 болт{]}} \\
            }
        }
    }
     \\ \\

    \theutterance \stepcounter{utterance}  
    & & & \multicolumn{2}{p{0.3\linewidth}}{
        \cellcolor[rgb]{0.9,0.9,0.9}{
            \makecell[{{p{\linewidth}}}]{
                \texttt{\tiny{[GM$|$GM]}}
                \texttt{proposal} \\
            }
        }
    }
    & & \\ \\

    \theutterance \stepcounter{utterance}  
    & & & \multicolumn{2}{p{0.3\linewidth}}{
        \cellcolor[rgb]{0.9,0.9,0.9}{
            \makecell[{{p{\linewidth}}}]{
                \texttt{\tiny{[GM$|$GM]}}
                [1, 0, 2, 0]
            }
        }
    }
    & & \\ \\

    \theutterance \stepcounter{utterance}  
    & & \multicolumn{4}{p{0.6\linewidth}}{
        \cellcolor[rgb]{0.9,0.9,0.9}{
            \makecell[{{p{\linewidth}}}]{
                \texttt{\tiny{[P1$\langle$GM]}}
                \texttt{Другой игрок отправил секретное предложение. Теперь ваша очередь отправить своё собственное предложение.} \\
            }
        }
    }
    & \\ \\

    \theutterance \stepcounter{utterance}  
    & \multicolumn{4}{p{0.6\linewidth}}{
        \cellcolor[rgb]{0.9,0.9,0.9}{
            \makecell[{{p{\linewidth}}}]{
                \texttt{\tiny{[P1$\rangle$GM]}}
                \texttt{{[}Предложение: 2 горшка, 1 сумка{]}} \\
            }
        }
    }
    & & \\ \\

    \theutterance \stepcounter{utterance}  
    & & & \multicolumn{2}{p{0.3\linewidth}}{
        \cellcolor[rgb]{0.9,0.9,0.9}{
            \makecell[{{p{\linewidth}}}]{
                \texttt{\tiny{[GM$|$GM]}}
                \texttt{proposal} \\
            }
        }
    }
    & & \\ \\

    \theutterance \stepcounter{utterance}  
    & & & \multicolumn{2}{p{0.3\linewidth}}{
        \cellcolor[rgb]{0.9,0.9,0.9}{
            \makecell[{{p{\linewidth}}}]{
                \texttt{\tiny{[GM$|$GM]}}
                [0, 1, 0, 2]
            }
        }
    }
    & & \\ \\

    \theutterance \stepcounter{utterance}  
    & & & \multicolumn{2}{p{0.3\linewidth}}{
        \cellcolor[rgb]{0.9,0.9,0.9}{
            \makecell[{{p{\linewidth}}}]{
                \texttt{\tiny{[GM$|$GM]}}
                [[0, 1, 0, 2], [1, 0, 2, 0]]
            }
        }
    }
    & & \\ \\

\end{supertabular}
}

\end{document}
